% Ad Atticum 11.17 --- headnote
% ad-atticum-11-17, 12 or 13 June 47 BC, Brundisium

\textit{Cicero to Atticus, written from Brundisium on
the day before the Ides or the Ides of June 47 BC ---
12 or 13 June (the manuscript dateline: \textit{Scr.\
Brundisi prid.\ Id.\ aut Id.\ Iun.\ a.\ 707 (47)}). A
single section, hurried into the hands of someone
else's couriers because they were on the point of
leaving, and made shorter still by Cicero's intention
to send his own messengers shortly. The occasion is
Tullia's arrival at Brundisium on the day before the
Ides (12 June): she has brought three letters from
Atticus and a long account of his kindness to her.}

\textit{What Tullia's presence draws out of her father
is not the pleasure such a daughter ought to give but
a grief sharper than her arrival warrants. The contrast
is set out in Cicero's most characteristic shape ---
her \textit{virtus}, \textit{humanitas}, and
\textit{pietas} ranged against \textit{tam misera
fortuna}, and the cause of that fortune assigned not
to any fault of hers but to the gravest fault of his
own. From this admission flows the letter's only
business: he expects neither consolation nor counsel
from Atticus any longer. He knows Atticus wants to
offer the one, and that no one can devise the other,
and that Atticus has tried everything, in many earlier
letters and most recently of all.}
